\subsection{Algebraically free families of extensions}

DEFINITION 6. --- Let L be an extension of a field K and \((\mathrm{E}_{i})_{i\in I}\) a family of subextensions of L. The family \((\mathrm{E}_{i})_{i\in I}\) is said to be algebraically free if the following condition is satisfied:

(AF) For each \(i \in I\) let \(\mathbf{A}_i\) be a subset of \(\mathbf{E}_i\) which is algebraically free over \(\mathbf{K}\) . Then \(\mathbf{A}_i \cap \mathbf{A}_j = \emptyset\) for \(i \neq j\) and \(\cup \mathbf{A}_i\) is algebraically free over \(\mathbf{K}\) .

Remark. --- By Prop. 3 (V, p.~107) it is enough to verify the condition (AF) for finite subsets \(A_{i}\) . We thus obtain the following result: if \((\mathrm{E}_{i})_{i\in I}\) is an algebraically free family, the same is true of \((\mathrm{E}_{i}^{\prime})_{i\in I}\) if \(E_{i}^{\prime}\) is a subextension of \(E_{i}\) for each \(i\in I\) ; conversely, if every family \((\mathrm{E}_{i}^{\prime})_{i\in I}\) , where \(E_{i}^{\prime}\) is a finitely generated subextension of E for each \(i\in I\) is algebraically free, then \((\mathrm{E}_{i})_{i\in I}\) is algebraically free. On the other hand, for \((\mathrm{E}_{i})_{i\in I}\) to be algebraically free it is necessary and sufficient that \((\mathrm{E}_{i})_{i\in J}\) should be algebraically free for every finite subset J of I. Speaking intuitively we may say that the algebraic independence of extensions is a property « of finite character ».

PROPOSITION 15. --- Let \((\mathbf{E}_i)_{i\in I}\) be a family of subextensions of a given extension L of a field K. Then the following conditions are equivalent:

\begin{enumerate}
\def\labelenumi{\alph{enumi})}
\item
  The family \((\mathrm{E}_i)_{i\in I}\) is algebraically free.
\item
  For each \(i \in I\) the extension \(E_i\) is algebraically disjoint over \(K\) from the extension \(F_i\) generated by the \(E_j\) for \(j \neq i\) .
\item
  There exists a family \((\mathbf{B}_i)_{i\in I}\) of disjoint subsets of \(L\) , such that \(\mathbf{B}_i\) is a transcendence basis of \(\mathbf{E}_i\) over \(K\) for each \(i\in I\) , and \(\mathbf{B} = \bigcup_{i\in I}\mathbf{B}_i\) is algebraically free over \(K\) .
\end{enumerate}

It is clear that \(a)\) implies \(c)\) .

Assuming \(c\) ), let us choose \(i\) in \(I\) ; put \(C_i = \bigcup_{j \neq i} B_j\) . For each \(j \neq i\) every element

of \(E_{j}\) is algebraic over \(\mathbf{K}(\mathbf{B}_{j})\) and a fortiori over \(\mathbf{K}(\mathbf{C}_{i})\) . By Cor. 1 of V, p.~18, the field \(F_{i}\) is therefore algebraic over \(\mathbf{K}(\mathbf{C}_{i})\) . Further, we have \(B_{i} \cap C_{i} = \varnothing\) and \(B = B_{i} \cup C_{i}\) is algebraically free over K; therefore \(B_{i}\) is algebraically free over \(\mathbf{K}(\mathbf{C}_{i})\) (V, p.~107, Prop. 4), hence also over \(F_{i}\) (which is algebraic over \(\mathbf{K}(\mathbf{C}_{i})\) ) by Prop. 6 of V, p.~108. Thus we have proved that \(E_{i}\) is algebraically disjoint from \(F_{i}\) over K (V, p.~113, Prop. 12), hence c) implies b).

Let us now assume \(b)\) and prove \(a)\) . It is enough to show that if \(i_1, \ldots, i_n\) are distinct elements of \(I\) , then the family of extensions \((E_{i_1}, \ldots, E_{i_n})\) is algebraically free; we argue by induction on \(n\) , the case \(n = 1\) being trivial. Suppose then that \(n > 1\) and that the family \((E_{i_1}, \ldots, E_{i_{n-1}})\) is algebraically free; for \(1 \leqslant k \leqslant n\) choose a subset \(A_k\) of \(E_{i_k}\) algebraically free over \(K\) and put

\(B = A_{1} \cup \ldots \cup A_{n-1}\) . By the induction hypothesis the subsets \(A_{1}, \ldots, A_{n-1}\) are pairwise disjoint and B is algebraically free over K; by b) \(E_{i_{n}}\) is algebraically disjoint from \(F_{i_{n}}\) and since B is contained in \(F_{i_{n}}\) we have \(B \cap A_{n} = \varnothing\) and \(B \cup A_{n} = A_{1} \cup \ldots \cup A_{n}\) is algebraically free over K. We have thus shown that the family \((E_{i_{1}}, \ldots, E_{i_{n}})\) is algebraically free.

The following proposition generalizes part \(b)\) of Prop. 13 (V, p.~113).

PROPOSITION 16. --- Let \((\mathrm{E}_i)_{i\in I}\) be a family of subextensions of an extension of a field K, and let E be the field generated by \(\cup\) \(\mathrm{E}_i\) .

\begin{enumerate}
\def\labelenumi{\alph{enumi})}
\item
  We have \(\operatorname{tr} \cdot \deg_{\mathbf{K}} \mathrm{E} \leqslant \sum_{i \in I} \operatorname{tr} \cdot \deg_{\mathbf{K}} \mathrm{E}_i\) , with equality when the family \((\mathrm{E}_i)_{i \in I}\) is algebraically free over \(\mathbf{K}\) .
\item
  Conversely, assume that \(\operatorname{tr} \cdot \deg_{\mathbf{K}} \mathrm{E} = \sum_{i \in I} \operatorname{tr} \cdot \deg_{\mathbf{K}} \mathrm{E}_i\) and that \(\operatorname{tr} \cdot \deg_{\mathbf{K}} \mathrm{E}\) is finite; then the family \((\mathrm{E}_i)_{i \in I}\) is algebraically free over \(\mathbf{K}\) .
\end{enumerate}

For each \(i \in I\) let \(\mathbf{B}_i\) be a transcendence basis of \(\mathbf{E}_i\) over \(\mathbf{K}\) and put \(\mathbf{B} = \bigcup_{i \in I} \mathbf{B}_i\) . For each \(i \in I\) , each element of \(\mathbf{E}_i\) is algebraic over \(\mathbf{K}(\mathbf{B}_i)\) , hence over

K(B). Now Cor. 1 of V, p.~18 shows that E is algebraic over K(B); by V, p.~109, Th. 2, B thus contains a transcendence basis of E over K. If moreover the family \((\mathrm{E}_{i})_{i\in I}\) is algebraically free over K, then the \(B_{i}\) are disjoint and the set B is algebraically free over K. This establishes a) (Set Theory, III, p.~160, Cor. of Prop. 4).

Under the hypotheses of \(b\) ), E is algebraic over \(\mathbf{K}(\mathbf{B})\) and of finite transcendence degree over \(\mathbf{K}\) , and we have \(\operatorname{Card}(\mathbf{B}) \leqslant \operatorname{tr} \cdot \deg_{\mathbf{K}} \operatorname{E}\) . By Cor. 1 of V, p.~110, B is algebraically free over \(\mathbf{K}\) and the \(\mathbf{B}_i\) are disjoint. Now Prop. 15 shows that the family \((\mathrm{E}_i)_{i \in I}\) is algebraically free over \(\mathbf{K}\) .

Before stating the next theorem let us remark that there exist algebraically closed extensions of K of arbitrary transcendence degree, for example an algebraic closure of an appropriate field of rational fractions.

THEOREM 5. --- Let \((\mathrm{E}_i)_{i\in I}\) be a family of extensions of a field \(\mathbf{K}\) and \(\Omega\) an algebraically closed extension of \(\mathbf{K}\) . Suppose that the inequality

\[
\operatorname{tr} \cdot \deg_ {K} \Omega \geqslant \sum_ {i \in I} \operatorname{tr} \cdot \deg_ {K} E _ {i}\tag{2}
\]

holds. Then there exists an algebraically free family \((\mathrm{F}_i)_{i\in I}\) of subextensions of \(\Omega\) such that \(\mathbf{F}_i\) is K-isomorphic to \(\mathbf{E}_i\) for all \(i\in I\) .

For each \(i \in I\) let \(B_i\) be a transcendence basis of \(E_i\) over \(K\) . Let \(B\) be a transcendence basis of \(\Omega\) over \(K\) . By (2) we have Card \(B \geqslant \sum_{i \in I} \text{Card } B_i\) ; hence

there exists a family \((\mathbf{B}_{i}^{\prime})_{i\in I}\) of pairwise disjoint subsets of B and bijections \(u_{i}:B_{i}\to B_{i}^{\prime}\) (for \(i\in I\) ). By Prop. 2 of V, p.~106, \(u_{i}\) extends to a K-isomorphism \(v_{i}\) of \(\mathbf{K}(\mathbf{B}_{i})\) onto \(\mathbf{K}(\mathbf{B}_{i}^{\prime})\) ; since \(\Omega\) is algebraically closed and \(E_{i}\) algebraic over

\(K(B_{i})\) , the Cor. (V, p.~23) shows that \(v_{i}\) extends to a K-isomorphism of \(E_{i}\) onto a subextension \(F_{i}\) of \(\Omega\) . By construction \(B_{i}^{\prime}\) is a transcendence basis of \(F_{i}\) over K, and Prop. 15 (V, p.~115) shows that the family \((F_{i})_{i\in I}\) of subextensions of \(\Omega\) is algebraically free over K.

COROLLARY 1. --- Let E and Ω be two extensions of a field K. Suppose that Ω is algebraically closed, of transcendence degree at least equal to that of E. Then E is K-isomorphic to a subextension of Ω.

COROLLARY 2. --- Let \(\Omega\) be an algebraically closed field of infinite transcendence degree over its prime subfield. Then every field of the same characteristic as \(\Omega\) is the ascending directed union of fields isomorphic to subfields of \(\Omega\) .

For each field is the ascending directed union of finitely generated subfields over its prime field, and now it is enough to apply Cor. 1.

* Example. --- This applies particularly in characteristic 0 in taking \(\Omega = \mathbf{C}\) (« Lefschetz' principle »). *

